% =====================================================================
\chapter{Methodology}
\label{ch:methodology}
% =====================================================================

This chapter presents the two methodological contributions of the thesis. First,
a \emph{per-component forgetting diagnostic} for tri-modal document encoders
(\S\ref{sec:method-diagnostic}). Second, a \emph{mechanism-targeted continual-learning
method}, DocCL, whose design is derived from the diagnosis rather than chosen in
advance (\S\ref{sec:method-remedy}). The ordering is deliberate: the diagnosis is
the primary contribution and constrains the remedy.

\section{Problem Formulation}
\label{sec:method-formulation}

We study continual \gls{kie} as token-level sequence labelling. It is useful to
separate three layers of ``input/output'' that are easily conflated: at the
\emph{user} level the input is a document scan and the output is extracted
information; at the \emph{machine-learning} level the input is a stream of
datasets $(\mathcal{D}_1,\mathcal{D}_2,\dots)$ and the output is the trained
parameters; at the \emph{technical} level---the one we formalise---the input is a
document tuple and the output is a label sequence. (All symbols introduced here
are collected in the List of Notations.)

A document is a triple $D = (W, B, I)$ where $W = (w_1,\dots,w_N)$ is a sequence of
$N$ \gls{ocr} tokens, $B = (b_1,\dots,b_N)$ with $b_i \in [0,1000]^4$ are
normalised bounding boxes, and $I \in \mathbb{R}^{3\times224\times224}$ is the page
image. The model $f_\theta$ produces \gls{bio} labels $Y = (y_1,\dots,y_N)$ with
$y_i \in \mathcal{L}_{\text{task}}$, trained by token-level cross-entropy $\ell$
with ignored positions masked out. Performance is measured by entity-level F1 with
\gls{bio}-aware span matching.

In the continual setting the model sees a sequence of tasks $\mathcal{T}_1, \dots,
\mathcal{T}_T$, each with its own data and (in \gls{cil}) its own growing label
space. Let $\theta$ denote the encoder parameters and $\theta^{(t)}$ the parameters
after training task $t$; the learner may access only $\mathcal{D}_t$ and (if any) a
memory $M$ at task $t$. The continual objective minimises the summed risk over all
tasks,
\[
  \theta_T = \arg\min_{\theta}\ \sum_{t=1}^{T}
    \mathbb{E}_{(x,y)\sim\mathcal{D}_t}\big[\ell(f_\theta(x), y)\big]
  \quad \text{s.t.}\quad |M| \le M_{\max},\ \ \lVert\Delta\theta\rVert \le P,
\]
under a bounded memory $M$ (at most $M_{\max}$ exemplars) and a parameter budget
$P$ (both relevant because the document setting favours memory-light,
privacy-preserving methods). The object of
study is how performance on earlier tasks degrades as later tasks are learned,
and---uniquely here---how that degradation is distributed across the
\emph{components} of $\theta$.

\subsection{Defining a Component}
\label{sec:method-components}

We partition the encoder parameters into groups that are \emph{genuinely
separable} in the backbone. A caveat governs this partition and is central to the
honesty of the analysis: LayoutLMv3 is a \emph{single-stream} encoder. Only the
\emph{input} embeddings are modality-specific; from the first encoder layer
onward, text and visual patch tokens are concatenated and share the \emph{same}
query/key/value projections. There is therefore no identifiable
``text-attention'' versus ``visual-attention'' versus ``cross-modal fusion''
\emph{parameter} group---those projections are one shared set. We do not pretend
otherwise. We group instead by what the architecture actually separates:

\begin{itemize}
  \item \textbf{Text word embeddings} --- the token embedding table
        (\texttt{text\_word\_embed}).
  \item \textbf{Layout embeddings} --- the 2D position embeddings encoding
        bounding-box geometry (\texttt{layout\_2d\_pos\_embed}).
  \item \textbf{Visual patch embedding} --- the image-patch projection
        (\texttt{image\_patch\_embed}).
  \item \textbf{Shared self-attention} --- the (single-stream) query/key/value
        projections (\texttt{attn\_qkv}) and the attention output projection
        (\texttt{attn\_out}), reported separately rather than folded into the
        feed-forward block.
  \item \textbf{Feed-forward and normalisation} --- the position-wise
        feed-forward blocks (\texttt{ffn}) and layer norms (\texttt{layernorm}).
  \item \textbf{Task head} --- the token-classification layer
        (\texttt{classifier}), tracked separately as it is expanded for new
        labels.
\end{itemize}

These groups are exhaustive and every one is populated; a residual catch-all is
verified empty for the stock backbone. The modalities' \emph{behavioural}
contributions are recovered separately, by the input-masking conditions of
\S\ref{sec:method-diag-ablation}, and forgetting is additionally localised by
\emph{depth} (input / early / mid / late encoder layers / head), the axis along
which a single-stream encoder is most naturally heterogeneous. For CKA tracking we
probe six depth-matched points (embeddings, patch embedding, early/mid/late
encoder layers, and the classifier) so the gradient can be compared across
backbones, including the unimodal BERT baseline.

\section{A Per-Component Forgetting Diagnostic}
\label{sec:method-diagnostic}

\subsection{Motivation}
\label{sec:method-diag-motivation}

If forgetting concentrates in a particular component, a continual-learning method
should spend its stability budget there; if it is uniform, uniform methods are
justified. Either way the answer is informative. To obtain it we need (i) a way to
isolate each modality's contribution without confounds, and (ii) metrics that
attribute forgetting to components.

\subsection{Controlled Modality Ablation}
\label{sec:method-diag-ablation}

We use two complementary controls. First, to isolate the modalities
\emph{behaviourally} without confounding architecture, we hold the architecture
fixed and mask the input streams of a \emph{single} LayoutLMv3, defining four
apples-to-apples conditions C1--C4 (Table~\ref{tab:diagnostic-conditions}): text
only (C1), image+layout (C2), text+layout (C3), and the full tri-modal input (C4).
Because only the inputs differ, any difference in forgetting is attributable to the
modalities in play rather than to the network. Critically, the text-only condition
(C1) keeps the real text tokens and zeroes \emph{only} the image and layout
streams---it is a genuine text-only LayoutLMv3, not a null input.

Second, because a same-architecture mask cannot tell us whether a multimodal
encoder forgets \emph{differently} from a unimodal one, we add an \emph{external}
unimodal baseline Cb: a true BERT-base text encoder, tokenised with its own
WordPiece vocabulary on the same documents and run through the same instrument
(per-token CKA, per-group importance, displacement, forgetting metrics). The
within-architecture conditions answer ``which stream/depth forgets''; the Cb
contrast answers ``is this multimodal-specific?''.

\begin{table}[!htbp]
  \centering
  \caption[The diagnostic conditions]{The diagnostic conditions. C1--C4 are the
  \emph{same} LayoutLMv3 architecture under different input-stream masks (only the
  active inputs differ); Cb is an external true-unimodal baseline. C1 keeps real
  text (image and layout zeroed)---it is not a null-input condition.}
  \label{tab:diagnostic-conditions}
  \begin{tabular}{cll}
    \toprule
    \textbf{ID} & \textbf{Model / active inputs} & \textbf{Purpose} \\
    \midrule
    Cb & BERT-base, text only          & external true-unimodal contrast \\
    C1 & LayoutLMv3, text only         & real text-only (image + layout zeroed) \\
    C2 & LayoutLMv3, image + layout    & isolate visual+layout-only forgetting \\
    C3 & LayoutLMv3, text + layout     & isolate text+layout-only forgetting \\
    C4 & LayoutLMv3, full tri-modal    & main subject of analysis \\
    \bottomrule
  \end{tabular}
\end{table}

Figure~\ref{fig:diagnostic-conditions} depicts the LayoutLMv3 conditions as the same
network with different input pathways switched on or off.

\begin{figure}[!htbp]
  \centering
  \includegraphics[width=\linewidth,height=0.5\textheight,keepaspectratio]{diagnostic_conditions.pdf}
  \caption[The diagnostic conditions as active and zeroed pathways]{The LayoutLMv3
  conditions C1--C4 realised on a single fixed architecture. Each panel shows the
  same LayoutLMv3 with a different subset of input pathways active (text, image,
  layout) and the remainder zeroed at the input: real text only (C1), image and
  layout (C2), text and layout (C3), and the full tri-modal input (C4). The external
  unimodal baseline Cb (BERT-base) is a separate model, not a mask. Because only the
  inputs differ across C1--C4, any difference in forgetting is attributable to the
  modalities in play rather than to the network.}
  \label{fig:diagnostic-conditions}
\end{figure}

\subsection{Forgetting Decomposition Metrics}
\label{sec:method-diag-metrics}

We measure forgetting with three complementary instruments, following the
multiple-instrument principle motivated in \S\ref{sec:bg-superficial-essential}.

\paragraph{Representational drift (per-token linear CKA).} For each tracked layer
we compute the linear CKA~\citep{kornblith2019similarity} between activations of
the model after task $t-1$ and after task $t$ on a fixed probe set. Given two
centred activation matrices $X \in \mathbb{R}^{n\times p}$ and $Y \in
\mathbb{R}^{n\times q}$,
\[
  \mathrm{CKA}(X,Y) = \frac{\lVert Y^\top X \rVert_F^2}
  {\lVert X^\top X \rVert_F \, \lVert Y^\top Y \rVert_F}.
\]
A drop of CKA below $1$ at a layer marks it as a drift hotspot. Two design choices
matter for validity. First, token classification lives in the \emph{per-token}
representation, so we compute CKA over \emph{individual valid tokens} (the
classifier-relevant first-subword positions), not over a document-mean vector that
would average $\sim$500 token vectors per page and could read ``no drift'' while
the token geometry has reorganised. Second, the probe collects $n\ge 500$ distinct
tokens (well above the hidden width $p=768$), so the estimate is not in the
high-variance $n<p$ regime; activations are captured \emph{under the condition's
modality mask} so each condition is probed on the inputs it was trained on.

\paragraph{Where forgetting lives: importance vs.\ displacement.} We separate two
distinct quantities that are easily conflated. The empirical Fisher
\emph{importance} of a group $g$ on a task's own data,
\[
  F_g = \frac{1}{|\mathcal{D}|}\sum_{(x,y)\in\mathcal{D}}
        \sum_{\phi\in g}\Big(\frac{\partial \log p_\theta(y\mid x)}{\partial \phi}\Big)^2 ,
\]
measures how \emph{sensitive} the loss is to $g$ \emph{for that task}. It is
\emph{not} a forgetting signal: the classifier head carries the largest such
gradients for \emph{any} token classifier (BERT included) by the geometry of the
cross-entropy output layer, regardless of whether those parameters were
\emph{moved}. To localise \emph{forgetting} we use the old-task-Fisher-weighted
\emph{displacement} between the checkpoint before and after a task,
\[
  D_g = \sum_{\phi\in g} F^{(t-1)}_\phi\,
        \big(\theta^{(t)}_\phi - \theta^{(t-1)}_\phi\big)^2 ,
\]
which is large precisely where parameters that mattered for the \emph{previous}
task ($F^{(t-1)}$) have been overwritten. We report $D_g$ per component group and
per depth bucket (input/early/mid/late/head) as the primary forgetting localiser,
with Fisher \emph{importance} reported alongside but never labelled forgetting.
Fisher is computed under each condition's modality mask.

\paragraph{Forgetting matrix.} We record the standard accuracy matrix
$R[i,j]$, the entity-F1 on task $j$ after training task $i$, from which we derive
average accuracy, backward transfer, and average forgetting
(\S\ref{sec:eval-metrics}). Its structure is annotated in
Figure~\ref{fig:forgetting-matrix}: the diagonal holds just-learned accuracy, the
lower triangle holds retained accuracy on earlier tasks (the source of backward
transfer and forgetting), and the upper triangle---unused here, as tasks arrive
sequentially---would hold zero-shot forward transfer. This is the conventional,
component-agnostic view that the per-component metrics refine.

\begin{figure}[!htbp]
  \centering
  \includegraphics[width=\linewidth,height=0.5\textheight,keepaspectratio]{forgetting_matrix_schematic.pdf}
  \caption[Annotated structure of the forgetting matrix]{Schematic structure of the
  accuracy matrix $R[i,j]$ (entity-F1 on task $j$ after training task $i$); symbolic
  labels only, with no measured values. The diagonal $R[i,i]$ records just-learned
  accuracy; the lower triangle $R[i,j]$ for $j<i$ records retained accuracy on
  earlier tasks, from which average accuracy, backward transfer, and average
  forgetting are derived; the upper triangle is unevaluated under sequential
  arrival.}
  \label{fig:forgetting-matrix}
\end{figure}

\subsection{Diagnostic Protocol}
\label{sec:method-diag-protocol}

For each condition C1--C4 and each seed we train naively and sequentially on
FUNSD\,$\rightarrow$\,CORD\,$\rightarrow$\,SROIE with no continual-learning
strategy, expanding the classifier head at each task. At every task boundary we
capture per-layer CKA, per-group Fisher, and the forgetting matrix.
Algorithm~\ref{alg:forgetting-decomposition} summarises the procedure and
Figure~\ref{fig:diagnostic-pipeline} shows it as a data-flow diagram.

\begin{algorithm}[t]
  \caption{Per-component forgetting decomposition}
  \label{alg:forgetting-decomposition}
  \begin{algorithmic}[1]
    \Require condition $c \in \{C1,C2,C3,C4\}$, task sequence $\mathcal{T}_{1:T}$,
             probe set $\mathcal{P}$, layer set $\mathcal{G}_{\text{CKA}}$,
             parameter groups $\mathcal{G}_{\text{F}}$
    \State initialise $\theta^{(0)}$ from the pre-trained backbone
    \For{$t = 1$ \textbf{to} $T$}
      \State expand classifier head to $\mathcal{L}_{\mathcal{T}_t}$
      \State $A_{t-1} \gets \textsc{Activations}(\theta^{(t-1)}, \mathcal{P})$
      \State $\theta^{(t)} \gets \textsc{FineTune}(\theta^{(t-1)}, \mathcal{T}_t)$
        \Comment{naive, modality mask $c$}
      \State $A_{t} \gets \textsc{Activations}(\theta^{(t)}, \mathcal{P})$
      \ForAll{$\ell \in \mathcal{G}_{\text{CKA}}$}
        \State $\mathrm{CKA}_{t,\ell} \gets \textsc{LinearCKA}(A_{t-1}[\ell], A_t[\ell])$
      \EndFor
      \ForAll{$g \in \mathcal{G}_{\text{F}}$}
        \State $F_{t,g} \gets \textsc{Fisher}(\theta^{(t)}, g, \mathcal{T}_t)$
      \EndFor
      \ForAll{$j \le t$}
        \State $R[t,j] \gets \textsc{EvalF1}(\theta^{(t)}, \mathcal{T}_j)$
      \EndFor
    \EndFor
    \State \Return $\{\mathrm{CKA}_{t,\ell}\}, \{F_{t,g}\}, R$
  \end{algorithmic}
\end{algorithm}

\begin{figure}[!htbp]
  \centering
  \includegraphics[width=\linewidth,height=0.5\textheight,keepaspectratio]{diagnostic_pipeline.pdf}
  \caption[The per-component diagnostic pipeline]{Data flow of the per-component
  diagnostic. The model is fine-tuned sequentially on
  FUNSD\,$\rightarrow$\,CORD\,$\rightarrow$\,SROIE under a fixed modality mask, and
  at each task boundary three instruments are captured: per-layer linear CKA between
  the pre- and post-task checkpoints on a fixed probe set, per-group empirical
  Fisher importance, and a row of the accuracy matrix $R[i,j]$ over all seen tasks.
  The procedure is repeated for every condition C1--C4 and every seed.}
  \label{fig:diagnostic-pipeline}
\end{figure}

\subsection{Hypotheses and Decision Rule}
\label{sec:method-diag-hypotheses}

We test a null hypothesis $H_0$ that forgetting is uniform across the network
against $H_1$ that it concentrates along depth (in particular at the classifier
head and late encoder layers). The test uses the per-group displacement $D_g$ and
per-depth displacement of \S\ref{sec:method-diag-metrics}, aggregated \emph{per
seed first} (the per-seed statistic is the mean over that seed's boundaries) so the
unit of analysis is the seed. For the dominant (largest mean) bucket we compare its
seed-level displacements against the pooled remaining buckets with a one-sided
Mann--Whitney U test, rejecting $H_0$ for a bucket when $p_g < 0.05/(K-1)$ for $K$
buckets.

\paragraph{Statistical power (a hard constraint on the design).} A rank test on
small samples has a \emph{minimum achievable} two-sided $p$-value of
$2/\binom{n_1+n_2}{n_2}$ even under perfect separation. With three seeds per group
this floor is $2/\binom{6}{3}=0.10$, far above any Bonferroni threshold: such a
test \emph{cannot} reject $H_0$ regardless of the data, so a ``failure to reject''
would be an artefact of sample size, not evidence. We therefore run \emph{at least
five seeds} per condition (and pool the alternate task order), which lowers the
floor for the cross-condition contrast to $2/\binom{10}{5}=0.0079$, below the
corrected threshold. Every reported non-rejection is accompanied by its power
floor, and any comparison whose floor exceeds $\alpha$ is labelled
\emph{inconclusive by construction} rather than evidence for $H_0$.

\paragraph{Location, not just magnitude.} Asking whether a multimodal encoder
forgets \emph{differently} is a question about the \emph{distribution} of
forgetting across components, which a scalar test on total backward-transfer
magnitude cannot answer. We therefore add a distribution-targeting test: for each
condition we form the per-depth displacement \emph{profile} (an $L_1$-normalised
vector over buckets) and compare the profile of C4 against each unimodal reference
(C1 and the external Cb) with a permutation test on the cosine distance between
mean profiles. This directly tests ``do they forget in a different \emph{place}?''.

The diagnostic outcome maps to the remedy through the measurable branches of
Table~\ref{tab:subhypotheses} and Figure~\ref{fig:selection-tree}. We deliberately
restrict the decision to signals the instrument can actually produce: because
LayoutLMv3 is single-stream, ``fusion-dominant'' or ``visual-attention-dominant''
branches are \emph{not realisable} and are excluded; the admissible outcomes are a
head/late-layer concentration, a uniform-across-depth pattern, or no concentration.

\begin{table}[!htbp]
  \centering
  \caption[Diagnostic outcomes and the remedy each selects]{Measurable diagnostic
  outcomes and the remedy each selects. Only outcomes the single-stream instrument
  can realise are admitted; the decision is a function of the pilot
  (\S\ref{sec:method-selection}).}
  \label{tab:subhypotheses}
  \begin{tabularx}{\textwidth}{l L l}
    \toprule
    \textbf{Pattern} & \textbf{Signature} & \textbf{Triggers} \\
    \midrule
    $H_a$ head/late-dominant & Displacement concentrates at the head and late layers & DocCL: depth/head-targeted (\S\ref{sec:method-selection}) \\
    $H_b$ uniform-across-depth & Displacement broadly even across depth buckets & Uniform consolidation \\
    $H_0$ no concentration & No bucket separates after correction (with adequate power) & Characterisation-only (\S\ref{sec:method-fallback}) \\
    \bottomrule
  \end{tabularx}
\end{table}

The mapping from diagnostic outcome to remedy is made explicit in
Figure~\ref{fig:selection-tree}: the decision is a function of the pilot, so the
mechanism is not committed to in advance, and the no-concentration branch leads to
the characterisation-only fallback rather than to a method.

\begin{figure}[!htbp]
  \centering
  \includegraphics[width=\linewidth,height=0.5\textheight,keepaspectratio]{method_selection_tree.pdf}
  \caption[Decision tree mapping diagnostic outcomes to the remedy]{The decision
  rule as a tree over \emph{measurable} outcomes. A head/late-layer concentration
  ($H_a$) selects the depth/head-targeted DocCL; a uniform-across-depth pattern
  ($H_b$) selects uniform consolidation; no significant concentration (at adequate
  power, $H_0$) selects the characterisation-only fallback. Unrealisable branches
  (fusion- or visual-attention-dominant) are excluded because a single-stream
  encoder cannot produce those signals. \emph{Figure to be regenerated from
  \texttt{figures\_src/drawio/} to match these branches.}}
  \label{fig:selection-tree}
\end{figure}

\section{Mechanism-Targeted Continual Learning}
\label{sec:method-remedy}

The diagnosis localises forgetting to a specific region of the network---the classifier
head and late encoder layers (Chapter~\ref{ch:results})---rather than prescribing a single
mechanism. That locus admits a \emph{family} of head/late-targeted remedies, and this thesis
develops two of them. The first, \textbf{DocCL} (\S\ref{sec:method-selected}), concentrates a
conventional consolidation/replay \emph{budget} on the locus. The second, \textbf{LexSlot}
(\S\ref{sec:method-lexslot}), protects the same locus with a parameter-efficient, isolated
\emph{slot memory} that needs no replay buffer. The two share the motivation, design-space
analysis, and selection criterion of \S\S\ref{sec:method-remedy-motivation}--\ref{sec:method-selection};
they differ only in how they spend their stability at the diagnosed locus, which makes the pair
a controlled test of \emph{how} to target, not merely \emph{where}.

\subsection{Why Component-Agnostic Continual Learning Is Insufficient}
\label{sec:method-remedy-motivation}

Generic continual-learning methods treat the encoder uniformly. EWC penalises all
important parameters with a single Fisher-weighted term regardless of which
component they belong to; LW-AFT freezes layers by a fixed heuristic. If
forgetting concentrates in one component, both spend stability budget where it is
not needed and may under-protect where it is. The hypothesis of DocCL is that
\emph{concentrating} the continual-learning mechanism---be it regularisation
strength, freezing, or replay weighting---on the dominant-forgetting component
improves the stability--plasticity trade-off at negligible extra cost.

\subsection{Design Space: Where a Mechanism Can Attach in a Single-Stream Encoder}
\label{sec:method-design-space}

The three continual-learning families of \S\ref{sec:bg-cl-families}---regularisation,
replay, and distillation---each leave one design choice open for a document encoder:
\emph{where} in the network the mechanism concentrates its effort. This is precisely
the choice that existing methods make blindly. EWC weights every parameter by its
Fisher importance but applies the \emph{same} quadratic form everywhere; replay
re-exposes whole inputs without regard to which layers most need grounding;
distillation matches the full output without asking which representations have
drifted. The contribution of this thesis is not a new family but a principled,
\emph{measured} answer to the ``where'' question---and the honest constraint on that
answer is the single-stream caveat of \S\ref{sec:method-components}.

That caveat rules out an entire class of otherwise-appealing designs. One might wish
to place per-modality low-rank adapters on a ``text'', ``visual'', and ``fusion''
bank, or to route per-modality prompt sub-pools by a layout-aware query---both natural
ideas, and both presupposing separable text-, visual-, and fusion-attention parameter
groups. LayoutLMv3 has none: from the first encoder layer the streams share one set of
query/key/value projections (\S\ref{sec:method-components}). We therefore do not
pursue modality-partitioned adapters or routers here; they are reserved for a
genuinely two-stream backbone where the streams physically exist
(Chapter~\ref{ch:conclusions}). What the architecture \emph{does} separate is
\emph{depth}---input embeddings, early/mid/late encoder blocks, and the task
head---and depth is the axis along which the diagnosis (\S\ref{sec:method-diagnostic})
can actually localise forgetting. The remedy is therefore built on the one axis the
instrument can measure.

Within that constraint the natural realisation of each family is a
\emph{component-specific} version of itself: a Fisher penalty whose strength varies by
depth bucket rather than a single global $\lambda$ (regularisation); a buffer that
re-exposes old classes at the drifting head rather than generic rehearsal (replay); and
an output distillation that protects the output-facing representations specifically
(distillation). These are exactly the three ingredients of the selected method
(\S\ref{sec:method-selected}). The mechanisms themselves are established; what is novel
is that \emph{where} each applies is fixed by measurement, not by a hand-tuned
heuristic or a uniform default.

\subsection{Testable Component Hypotheses}
\label{sec:method-component-hypotheses}

Two method-level hypotheses are pre-registered and map one-to-one to the
component-targeting ablation in Chapter~\ref{ch:results}:

\begin{itemize}
  \item \textbf{H\textsubscript{target}.} Concentrating the stability budget on the
        diagnosed locus (the classifier head and late encoder layers) outperforms
        spreading the penalty uniformly across depth (the \emph{uniform}
        ablation, which penalises every bucket equally) by $\ge 2$ entity-F1 on
        $\ge 2/3$ scenarios.
  \item \textbf{H\textsubscript{head}.} Head-only protection recovers most of the
        gain of the full method; adding late-layer protection and head replay adds
        a further measurable improvement (the \emph{head-only} vs.\ \emph{late-only}
        vs.\ \emph{full} ablation).
\end{itemize}

A falsified hypothesis remains informative and is reported as such.

\subsection{Selection Criterion}
\label{sec:method-selection}

The remedy follows directly from the measurable diagnostic outcome
(Table~\ref{tab:subhypotheses}): a head/late-layer concentration selects the
depth/head-targeted DocCL of \S\ref{sec:method-selected}; a uniform-across-depth
pattern selects uniform consolidation; no concentration (at adequate power) selects
the characterisation-only fallback (\S\ref{sec:method-fallback}). Beyond alignment
with the diagnosis, the method must be implementable within the compute budget,
defensible under reviewer scrutiny, and reuse the existing EWC/replay/distillation
infrastructure.

\subsection{The Selected Method: DocCL}
\label{sec:method-selected}

The diagnosis localises forgetting in the \emph{classifier head and late encoder
layers}, with the input encoders stable (Chapter~\ref{ch:results}). DocCL therefore
spends its stability budget there. Its three ingredients are deliberately
\emph{established} mechanisms---elastic-weight consolidation, output distillation, and
replay; the contribution is not a new mechanism but the diagnosis-driven decision of
\emph{where} each applies (the head and late layers, fixed by measurement rather than
by a uniform default or a hand-tuned freezing heuristic), together with the
component-targeting ablation that isolates the effect of that targeting. The three
mechanisms are mutually-reinforcing and each individually ablatable:

\begin{enumerate}
  \item \textbf{Depth-scaled elastic-weight consolidation.} A Fisher-weighted
        quadratic anchor to the previous task whose per-bucket weight scales with
        where forgetting lives---high on the head and late layers, near zero on the
        stable input encoders---using the depth grouping of
        \S\ref{sec:method-components}.
  \item \textbf{Output distillation.} A temperature-softened KL term matching the
        previous task's output distribution from a frozen teacher, which protects
        the output-facing (head and late) representations functionally.
  \item \textbf{Head re-exposure replay.} A small reservoir buffer---replay being
        the strongest baseline in the diagnosis---keeps the head grounded in earlier
        tasks.
\end{enumerate}

The total objective is $\mathcal{L} = \mathcal{L}_{\mathrm{CE}}
+ \tfrac{\lambda}{2}\sum_g w_g\, \mathcal{L}^{\mathrm{EWC}}_g
+ \alpha\,\mathcal{L}_{\mathrm{KD}} + \beta\,\mathcal{L}_{\mathrm{replay}}$,
where $\mathcal{L}_{\mathrm{replay}}$ is the cross-entropy on a small reservoir buffer
$\mathcal{B}$ (weight $\beta$, buffer size $|\mathcal{B}|$; values in
Appendix~\ref{app:hyperparameters}) and the depth weights $w_g$ are concentrated on the
diagnosed locus: with the depth buckets of \S\ref{sec:method-components},
\[
  w = \{\,\text{input}:0,\ \text{early}:0,\ \text{mid}:0.5,\ \text{late}:1.0,\ \text{head}:2.0\,\},
\]
so the head is anchored twice as strongly as the late layers and the stable input
encoders are left free to adapt. A \texttt{target\_depth} control (\emph{full},
\emph{head-only}, \emph{late-only}, \emph{uniform}) sets these weights and drives the
component-targeting ablation that tests H\textsubscript{target} and
H\textsubscript{head}: the \emph{uniform} setting applies an \emph{equal} non-zero
penalty to every depth bucket---including the input and early encoders that the
targeted schedule leaves free---rather than concentrating the budget on the head/late
locus. It thus spends, if anything, a \emph{larger} total penalty than the targeted
schedule ($\sum_g w_g = 5.0$ uniform versus $3.5$ targeted), so a collapse under
uniform isolates \emph{where} the budget is spent, not how much, as the operative
factor; the mechanisms themselves are held fixed across settings. The
method reuses the EWC Fisher machinery, the LwF-style teacher, and the shared
reservoir buffer, so it adds no fragile new infrastructure.

\subsection{Pre-Registered Fallback (Not Triggered)}
\label{sec:method-fallback}

For completeness we record the contingency registered before the diagnosis was run:
had forgetting \emph{not} concentrated ($H_0$ not rejected at adequate power), the
diagnostic of \S\ref{sec:method-diagnostic} and the FCS-CL benchmark would have become
the headline contributions and the remedy would have been withheld. The diagnosis did
concentrate forgetting at the head and late layers (Chapter~\ref{ch:results}), so this
branch was \emph{not} taken: DocCL is the selected, depth/head-targeted remedy. We
state the fallback explicitly because pre-registering it---rather than choosing the
method after seeing which design won---is what licenses the claim that DocCL's form was
\emph{derived from} the measurement; and because the diagnostic's value stands
independently of the remedy's margin regardless.

\section{Lexical-Slot Head-Memory Continual Learning (LexSlot)}
\label{sec:method-lexslot}

DocCL protects the diagnosed locus by spending a consolidation and replay \emph{budget}
there: a depth-scaled Fisher anchor, a distillation teacher, and a small rehearsal buffer.
That budget is effective but not free---it keeps a frozen teacher in memory and rehearses
stored inputs. LexSlot asks a complementary question: \emph{can the same head/late locus be
protected by a parameter-efficient structured memory that stores no past inputs at all?} It
shares the diagnosis, the design-space constraint, and the selection criterion of DocCL
(\S\S\ref{sec:method-remedy-motivation}--\ref{sec:method-selection}); it differs only in the
mechanism it attaches at the locus---an \emph{isolated, write-once slot memory} in place of a
consolidation budget.

\subsection{Slot Memories at the Diagnosed Locus}
\label{sec:method-lexslot-slots}

LexSlot freezes the backbone and inserts small additive \emph{slot} modules exactly where the
diagnosis places forgetting. Two kinds of slot are used. \textbf{Head logit-slots} attach to the
classifier: a slot bank of $n_{\text{head}}$ learnable value vectors produces a per-token logit
shift $\Delta\ell = \big(f V^\top\big)P$ from the per-token features $f$ entering the head
(value matrix $V$, projection $P$), added to the head's output. \textbf{Late-layer
representation-slots} attach to each encoder block in the diagnosed \emph{late} third: a
rank-$r$ adapter adds $\Delta h = (h\,D)\,U$ to the block's hidden state ($D\in\mathbb{R}^{d\times
r}$ down, $U\in\mathbb{R}^{r\times d}$ up). All slot up-projections are \emph{zero-initialised}
(LoRA-style), so an untrained slot is an exact no-op and adding slots never perturbs the frozen
network until they are trained. Placement is governed by a \texttt{slot\_depth} control
(\emph{head\_only}, \emph{head\_late}, \emph{uniform}); the canonical setting is
\emph{head\_late} (head $+$ late third), matching DocCL's diagnosed target, and \emph{uniform}
(slots on every layer) is the targeting ablation.

\subsection{Per-Task Isolation: Write-Once Slots}
\label{sec:method-lexslot-isolation}

Each task \emph{owns} a fixed, disjoint budget of $n_{\text{slots}}/T$ slots, claimed from the
unclaimed pool when the task begins ($T$ the expected number of tasks). In the selected
configuration each task updates \emph{only its own} slots: a per-slot gradient mask freezes every
other task's slots during training, so a new task can never overwrite the parameters a previous
task wrote. Because the backbone is frozen and prior slots are frozen, earlier tasks are
\emph{structurally} protected---there is no consolidation penalty to tune and no replay buffer to
store. At inference every slot contributes additively (no routing or task-id is required), so a
document is processed by the union of all tasks' slots over the shared frozen backbone. This is
the sense in which LexSlot is \emph{write-once}: knowledge for a task is written into its own
slots once and never disturbed.

\subsection{The Lexical Gate and the Sharing Ablation}
\label{sec:method-lexslot-gate}

Isolation is one extreme of a more general gate. LexSlot computes, for each task, a sparse
\emph{lexical signature}---a TF-IDF bag-of-tokens over the task's OCR input tokens---and the
cosine similarity $S_{ij}$ between task signatures. This signature is purely input-derived and
therefore \emph{drift-immune}: it does not depend on any learned, forgettable parameter. The
\texttt{slot\_sharing} control then sets the gradient mask on a prior task's slots when training a
new one: \emph{off} freezes them (the isolated, selected method); \emph{soft} makes them trainable
in proportion to $S_{ij}$ (lexically similar tasks co-train shared slots); \emph{hard} unlocks them
above a similarity threshold. The motivation for graded sharing is that on this benchmark the
lexical overlap is itself label-bearing---the receipt tasks SROIE and CORD share the
\texttt{total}/\texttt{tax}/\texttt{cash} vocabulary ($S\!\approx\!0.31$) while the form task FUNSD
is an outlier ($S\!\approx\!0.10$)---so high-$S$ tasks might \emph{help} each other by refining a
shared lexical-to-label rule. Whether that putative transfer outweighs the interference of letting
two tasks write the same slots is an empirical question, and the \texttt{slot\_sharing} axis is the
controlled ablation that answers it (Chapter~\ref{ch:results}). We present \emph{off} (isolation) as
the method because it is what the measurement selects; \emph{soft} and \emph{hard} are reported as
the ablation that justifies the choice.

\subsection{Relation to DocCL}
\label{sec:method-lexslot-vs-doccl}

LexSlot reuses DocCL's training scaffolding---the same forward pass and the same optional
consolidation terms are available---but its operative mechanism is the slot memory, and in the
selected configuration the protection of old tasks comes entirely from slot isolation rather than
from a Fisher penalty or rehearsal. The contrast is deliberate and ablatable: DocCL answers
``concentrate a consolidation/replay budget at the locus'', LexSlot answers ``install an isolated
parameter-efficient memory at the locus''. Both are diagnosis-derived---neither would be motivated
without the head/late localisation---so reporting them together tests whether the diagnosis admits
more than one effective realisation. Slot counts, ranks, and the per-task budget are listed in
Appendix~\ref{app:hyperparameters}; the implementation is released alongside DocCL.

\section{Summary}
\label{sec:method-summary}

We have defined the genuinely separable components of a single-stream tri-modal
document encoder, a controlled modality-ablation protocol with an external unimodal
contrast, and complementary metrics---per-token CKA, per-group importance, an
old-task-Fisher-weighted displacement (the forgetting localiser), and the
forgetting matrix---that together localise forgetting by depth, stream, and head.
The diagnosis admits a family of head/late-targeted remedies, of which we develop two: DocCL,
which concentrates a consolidation/replay stability budget on the diagnosed locus under a
pre-registered, adequately powered decision rule; and LexSlot, which protects the same locus with
an isolated, buffer-free slot memory gated by a drift-immune lexical signature. Reporting the pair
tests whether the diagnosis admits more than one effective realisation. The next chapter describes
how this methodology is realised in software.
